% Legacy core manuscript fragment retained only for provenance.
% It assumes the theorem environments and notation commands defined in main.tex.
% Replace the corresponding old sections before including this file, because several
% labels intentionally retain their old names for cross-reference compatibility.

\begin{abstract}
Selective segmentation requires more than an accurate mask: a system must also rank
its predictions by the harm they would incur if accepted. That ranking depends on the
deployment loss. A confidence useful for Dice need not be useful for a boundary loss,
and an area-based confidence is not automatically meaningful when the application
cares about surface error. We formalize this dependence by indexing both selective
risk and confidence by a bounded, measurable segmentation loss $L$. The ideal score
is the negative conditional risk of the mask actually deployed,
$C_L^\star(x)=-\mathbb{E}[L(Y,\widehat Y(x))\mid X=x]$.

The true label posterior needed by this definition is unavailable from a single
probability map. We therefore introduce the level-set law as an explicit
\emph{working posterior}: draw one threshold $T\sim\mathrm{Unif}(0,1)$ and set
$Y_T=\{i:p_i\geq T\}$. Under this stated model, expected loss is exactly the
one-dimensional integral $\int_0^1 L(Y_t,\widehat Y)\,dt$, which yields a common
$M$-point quadrature estimator for any such loss. We separate its two sources of
error: a deterministic quadrature term, bounded by $V(H_x)/(2M)$ for midpoint
quadrature, and a loss-specific discrepancy between the working and true posteriors.
Thus the numerical approximation can be controlled without claiming that the
working posterior is the true one.

For Dice, the framework estimates an expectation of a ratio. Soft Dice Confidence
(SDC), its closest predecessor, is instead a ratio of expectations. We give their
exact relationship, a finite counterexample showing that they generally differ, and
clarify that their probabilistic assumptions are not ordered: our construction is
more general in the deployment loss, but not uniformly weaker in its assumptions.
The result is a loss-indexed formulation of post-hoc, image-level confidence that
keeps the deployed decision, the evaluation risk, and the confidence score explicit.
\end{abstract}

\section{Introduction}

A segmentation system deployed with human oversight must decide which predictions
can be accepted and which should be deferred. Average test accuracy does not answer
this question. Two systems with the same average Dice can have very different value
if one reliably sends its failures to a human and the other assigns those failures
high confidence. Selective prediction captures this requirement: a fixed model emits
a prediction $\widehat Y(x)$ and an image-level confidence $g(x)$; predictions are
accepted from high to low confidence, and the induced ranking is evaluated through a
risk--coverage curve \cite{selectivenet,feng,jaeger}.

The risk in that curve is application-dependent. In segmentation, Dice and IoU
measure overlap, whereas surface Dice, average symmetric surface distance, and
HD95 measure boundary behavior. The distinction is operational rather than
terminological. Moving a small component far from the target can change little in
overlap while producing a large surface error; a spatially diffuse boundary error can
have the opposite profile. Consequently, a confidence score cannot be called
``correct'' without specifying the loss under which accepted predictions are judged.

Most single-pass post-hoc scores aggregate a per-pixel quantity such as maximum
probability or entropy. Systematic studies have shown that aggregation is itself a
major design choice \cite{values,zenk}. Soft Dice Confidence (SDC) takes a more direct
route: it aligns confidence with the Dice objective and provides an efficient
Dice-specific marginal surrogate \cite{sdc}. Its success motivates a broader
question:

\begin{quote}
Can the same metric-alignment principle produce a post-hoc confidence for an
arbitrary bounded segmentation loss, while making clear what posterior assumptions are
required?
\end{quote}

We answer this question by separating three objects that are often conflated.
First, the \emph{deployed action} $\widehat Y(x)$ is the mask that will actually be
accepted or rejected. It may be obtained with a threshold of $0.5$, a threshold of
$0.3$, an argmax, or a fixed post-processing pipeline; the theory assigns no special
status to one decision rule. Second, the \emph{deployment loss} $L$ specifies what a
bad accepted prediction means. Third, a distribution over possible labels is needed
to evaluate the conditional risk of the deployed action before its label is observed.

The first two objects are specified by the task. The third is not supplied by a
standard segmentation network: a probability map gives one number per pixel but not
the joint distribution of an entire mask. We use a classical random-set construction
as an explicit working model. A single uniform threshold applied to the probability
map generates a nested family of candidate masks. This construction preserves the
map as its one-pixel coverage function and turns expected loss into a one-dimensional
threshold integral \cite{goodman,molchanov}. We do not claim that the resulting law
is the true label posterior. Its role is analogous to that of a working likelihood:
it makes the assumed dependence structure visible and produces a computable target.

This perspective gives a simple loss-indexed estimator. For selected nodes $t_m$ and
weights $w_m$,
\begin{equation*}
  C_{L,M}(x)
  =-\sum_{m=1}^M w_m L\!\left(Y_{t_m},\widehat Y(x)\right).
\end{equation*}
Only $L$ changes between an overlap confidence and a boundary confidence. Every
candidate label is compared with the same deployed mask; pairwise dispersion among
candidate labels is a different estimand.

The formulation also prevents an overly strong interpretation of numerical
convergence. Increasing $M$ improves the approximation of expected loss under the
working posterior. It cannot repair a miscalibrated probability map or an incorrect
joint model. We make this separation explicit with a total error bound containing a
quadrature term and a loss-specific posterior discrepancy. The latter is deliberately
weaker than requiring the two complete posteriors to be close in total variation: for
selective prediction, only their expected loss for the action actually deployed is
directly relevant.

Our contributions are:
\begin{enumerate}
\item We define loss-indexed selective risk $\AURC_L$ and its Bayes-optimal
image-level confidence $C_L^\star$. This makes the deployment loss and deployed mask
part of the problem statement rather than implicit evaluation choices.
\item We introduce a unified level-set risk estimator for any bounded, measurable
segmentation loss with an explicit empty-set convention. Under the level-set working
posterior its target is exact, and midpoint quadrature has a deterministic
$O(1/M)$ error bound. A separate discrepancy term states precisely what must be
assumed for proximity to the true ideal confidence.
\item We establish the exact relationship to SDC. SDC is a ratio of expected Dice
components and therefore depends only on pixel marginals; the level-set Dice score is
an expected Dice and depends on a joint law. Neither estimator's probabilistic
assumptions imply the other's.
\end{enumerate}

The contribution is therefore not a claim that one posterior model is universally
correct, nor that one confidence must dominate for every risk. It is a framework for
constructing and testing a sharper claim: a confidence should be indexed by the loss
used to evaluate selective decisions.

\section{Loss-indexed selective segmentation}\label{sec:setup}

\paragraph{Prediction and loss.}
Let $(X,Y)$ be an image and its binary segmentation mask, drawn from an unknown
distribution. A fixed segmentation system returns a probability map
$p(X)\in[0,1]^N$ and a deployed hard prediction $\Yhat(X)$. For example,
$\Yhat(X)=\{i:p_i(X)\geq\gamma\}$ for a fixed binary decision threshold $\gamma$,
but all definitions below apply to any deterministic deployed decision rule.

Let $L_x(A,B)$ be the loss incurred when $B$ is deployed and $A$ is the label. We
assume that $L_x$ is measurable and bounded,
\begin{equation}\label{eq:bounded-assumption}
  0\leq L_x(A,B)\leq D_x<\infty.
\end{equation}
We suppress the subscript $x$ when no ambiguity arises. Boundedness is automatic for
$1-\Dice$ and $1-\IoU$. A surface loss that is undefined for an empty mask must be
extended before it enters the framework; for example, one can assign a one-sided
empty comparison the image diagonal. The choice is part of the definition of $L$ and
must be shared by the confidence computation and the selective-risk evaluation.

\paragraph{Risk and coverage are loss-indexed.}
A confidence $g:\mathcal X\to\mathbb R$ is oriented so that larger values mean more
confident predictions. At threshold $\tau$, define
\begin{equation}\label{eq:risk-coverage}
\begin{split}
  \mathrm{cov}_g(\tau)
  &=\Prb\!\left(g(X)\geq\tau\right),\\
  R_L(\tau;g)
  &=\frac{\Ex\!\left[
      L(Y,\Yhat(X))\mathbf 1\{g(X)\geq\tau\}
    \right]}
    {\Prb(g(X)\geq\tau)} .
\end{split}
\end{equation}
Randomization within a tied confidence level can be used to attain any desired
coverage. Writing $R_L(c;g)$ for the resulting risk at coverage $c$, the
loss-indexed area under the risk--coverage curve is
\begin{equation}\label{eq:aurc-l}
  \AURC_L(g)=\int_0^1 R_L(c;g)\,dc .
\end{equation}
% TODO(citation): add a canonical primary reference for the risk--coverage/AURC
% definition used in selective prediction.
The subscript is essential: $\AURC_{1-\Dice}$ and $\AURC_{\HDF}$ evaluate different
orderings of harm and have different units. Raw values should not be averaged across
heterogeneous losses without an explicit normalization.

In a finite evaluation set, sorting images by $g$ and averaging the loss over the top
$k$ predictions gives the empirical risk at coverage $k/n$. Consequently,
\begin{equation}\label{eq:rank}
  \AURC_L(g)\text{ depends on }g\text{ only through its induced ranking,}
\end{equation}
up to the stated tie rule. Any strictly increasing transformation of $g$ gives the
same empirical AURC.

\paragraph{Ideal confidence.}
Before observing the label, the relevant quantity is the conditional risk of the
mask actually deployed.
\begin{definition}[Ideal loss-indexed confidence]\label{def:ideal}
For a fixed deployed predictor $\Yhat$, define
\begin{equation}\label{eq:ideal}
  \rho_L^\star(x)
  :=\Ex\!\left[L(Y,\Yhat(x))\mid X=x\right],
  \qquad
  C_L^\star(x):=-\rho_L^\star(x).
\end{equation}
\end{definition}
If the deployment threshold changes from $\gamma=0.5$ to $\gamma=0.3$, then
$\Yhat$ changes and so does $C_L^\star$. The confidence is decision-aware; it is not
only a property of the uncertainty distribution.

\begin{proposition}[Optimal selective ordering]\label{prop:ideal-ranking}
Among all acceptance rules measurable with respect to $X$, accepting images in
increasing order of $\rho_L^\star(X)$ minimizes selective risk at every coverage,
up to randomization among ties. Equivalently, ranking by $C_L^\star$ produces a
pointwise optimal risk--coverage curve and hence minimizes $\AURC_L$.
\end{proposition}
\begin{proof}
For an acceptance function $a(X)\in[0,1]$ with $\Ex[a(X)]=c$, conditioning on $X$
gives
\begin{equation*}
  \Ex\!\left[L(Y,\Yhat(X))a(X)\right]
  =\Ex\!\left[\rho_L^\star(X)a(X)\right].
\end{equation*}
If an accepted point has larger conditional risk than a rejected point, exchanging
an equal amount of acceptance probability between them weakly decreases the
objective while preserving coverage. Repeating this exchange yields a threshold on
$\rho_L^\star$, with randomization only at a tied threshold. Division by the fixed
coverage $c$ proves the risk statement; integration over $c$ proves the AURC
statement.
\end{proof}

Definition~\ref{def:ideal} is a target, not an available score. It requires the full
conditional law of $Y$ given $x$. A standard probability map provides at most
estimates of the pixel marginals
$q_i(x)=\Prb(Y_i=1\mid X=x)$; those marginals do not identify how pixels vary jointly,
and nonlinear losses such as Dice and HD95 depend on that joint structure. The next
section makes one working choice explicit.

\section{Level-set risk framework}\label{sec:framework}

\subsection{An explicit working posterior}\label{sec:posterior}

For a fixed image, write $p=(p_1,\ldots,p_N)$. We use a shared random threshold to
construct a law over masks.

\begin{definition}[Level-set working posterior $Q_p$]\label{def:levelset}
Let $T\sim\mathrm{Unif}(0,1)$ and define
\begin{equation}\label{eq:levelset}
  Y_T=\{i:p_i\geq T\},
  \qquad
  Y_t=\{i:p_i\geq t\}.
\end{equation}
The distribution of $Y_T$ is denoted by $Q_p$.
\end{definition}

This construction has two immediate properties:
\begin{equation}\label{eq:levelset-marginals}
  \Prb_{Q_p}(i\in Y_T)=p_i,
  \qquad
  \Prb_{Q_p}(i,j\in Y_T)=\min(p_i,p_j).
\end{equation}
Thus $p$ is exactly the one-pixel coverage function of $Q_p$, while the common
threshold selects the comonotone coupling of those marginals. Samples are nested:
if $s<t$, then $Y_t\subseteq Y_s$. The object is a classical coverage-function
random set \cite{goodman,molchanov}; related shared-level constructions are used for
uncertain contours \cite{poethkow,bolin}.

Equation~\eqref{eq:levelset-marginals} is an identity internal to $Q_p$, not evidence
that $p_i$ equals the true marginal $q_i$ or that the true labels are comonotone.
Calling $Q_p$ a working posterior is meant to keep this distinction visible. It is
attractive because it generates spatially coherent, nested candidates at negligible
sampling cost. It can be inaccurate when real ambiguity changes topology, moves
several components independently, or is poorly represented by the probability map.

\subsection{Expected loss as a threshold integral}

\begin{proposition}[Threshold-integral risk]\label{prop:integral}
For any loss satisfying \eqref{eq:bounded-assumption},
\begin{equation}\label{eq:integral}
  r_{L,Q}(x)
  :=\Ex_{Y\sim Q_{p(x)}}\!\left[L(Y,\Yhat(x))\right]
  =\int_0^1 L(Y_t,\Yhat(x))\,dt .
\end{equation}
\end{proposition}
\begin{proof}
By Definition~\ref{def:levelset}, a draw $Y\sim Q_p$ is $Y_T$ with
$T\sim\mathrm{Unif}(0,1)$. Taking the expectation over $T$ gives the integral.
Boundedness implies integrability.
\end{proof}

The proposition is exact for the stated working posterior and does not depend on a
particular algebraic form of $L$. For example, one may insert $1-\Dice$, $1-\IoU$,
a bounded surface distance, or a task-specific asymmetric cost. This is the sense in
which the framework is loss-general. It does not say that the integral equals the
true conditional risk.

Let
\begin{equation*}
  H_{L,x}(t):=L(Y_t,\Yhat(x)).
\end{equation*}
For nodes $t_m\in(0,1)$ and nonnegative weights $w_m$ with
$\sum_m w_m=1$, define
\begin{equation}\label{eq:quad}
  r_{L,M}(x)=\sum_{m=1}^M w_m H_{L,x}(t_m),
  \qquad
  C_{L,M}(x)=-r_{L,M}(x).
\end{equation}
Every candidate mask is compared with the fixed deployed prediction $\Yhat(x)$.
By contrast, averaging $L(Y_{t_m},Y_{t_n})$ over pairs of level sets estimates
posterior dispersion; it does not estimate the conditional risk of the deployed
action in Definition~\ref{def:ideal}.

A simple parameter-free numerical rule is equal-width midpoint quadrature,
\begin{equation}\label{eq:midpoint-nodes}
  t_m=\frac{m-1/2}{M},
  \qquad
  w_m=\frac1M,
  \qquad m=1,\ldots,M.
\end{equation}
Here ``parameter-free'' means that the nodes are fixed by the numerical rule once
$M$ is chosen; it does not mean that midpoint nodes are optimal for AURC or for every
finite-$M$ integrand.

For a finite image, $H_{L,x}$ is a bounded step function: its value changes only
when $t$ crosses a distinct value in $p(x)$. It therefore has finite total variation
$V(H_{L,x})$. The one-dimensional Koksma bound for the midpoint rule gives
\begin{equation}\label{eq:koksma}
  \left|r_{L,M}(x)-r_{L,Q}(x)\right|
  \leq \frac{V(H_{L,x})}{2M}.
\end{equation}
% TODO(citation): add a primary reference for the one-dimensional
% Koksma--Hlawka inequality and midpoint star discrepancy 1/(2M).
This proves convergence to the working-posterior risk as $M$ grows. It does not
provide a uniform numerical guarantee unless $V(H_{L,x})$ is itself controlled, and
it does not imply that empirical AURC improves monotonically with $M$.

\subsection{Numerical error and posterior discrepancy}\label{sec:error}

Let $P_x=\Prb(Y\in\cdot\mid X=x)$ denote the true label posterior and define the
loss-specific discrepancy for the deployed action
\begin{equation}\label{eq:loss-discrepancy}
  \Delta_{L,\Yhat}(x)
  :=\left|
  \Ex_{Y\sim P_x}L(Y,\Yhat(x))
  -\Ex_{Y\sim Q_{p(x)}}L(Y,\Yhat(x))
  \right|.
\end{equation}
This quantity asks only whether the two posteriors agree on the expected loss of the
action being scored. It is weaker than requiring the full distributions to be close.
For example, if
$\mathrm{TV}(P_x,Q_{p(x)})\leq\delta_x$, with total variation defined as the
supremum over events, then boundedness gives the sufficient but potentially loose
bound
\begin{equation}\label{eq:tv-sufficient}
  \Delta_{L,\Yhat}(x)\leq D_x\delta_x.
\end{equation}
The converse need not hold: two posteriors can disagree strongly on events irrelevant
to the chosen loss and action while having zero discrepancy in
\eqref{eq:loss-discrepancy}.

\begin{proposition}[Conditional total-error bound]\label{prop:total-bound}
For midpoint quadrature and any bounded measurable $L$,
\begin{equation}\label{eq:decomp}
  \left|C_{L,M}(x)-C_L^\star(x)\right|
  \leq
  \frac{V(H_{L,x})}{2M}
  +\Delta_{L,\Yhat}(x).
\end{equation}
Consequently, under the stronger sufficient condition
$\mathrm{TV}(P_x,Q_{p(x)})\leq\delta_x$, the second term may be replaced by
$D_x\delta_x$.
\end{proposition}
\begin{proof}
Negating risks does not change their absolute difference. Insert the
working-posterior risk $r_{L,Q}(x)$ between $r_{L,M}(x)$ and the true conditional
risk, apply the triangle inequality, then use \eqref{eq:koksma} and
\eqref{eq:loss-discrepancy}. Equation~\eqref{eq:tv-sufficient} gives the final
statement.
\end{proof}

The bound exposes an irreducible distinction. Increasing $M$ can make the first term
arbitrarily small. It does not change $p$, the coupling selected by $Q_p$, or the
second term. Perfect marginal calibration alone also does not make the second term
zero for a non-additive loss: multiple joint posteriors can share the same calibrated
marginals and assign different probabilities to whole masks.

For AURC, numerical accuracy matters through ordering. If an available bound
$\varepsilon_i$ satisfies
$|r_{L,M}(x_i)-\rho_L^\star(x_i)|\leq\varepsilon_i$, then the order of two images is
certified whenever
\begin{equation}\label{eq:ranking-margin}
  r_{L,M}(x_j)-r_{L,M}(x_i)>\varepsilon_i+\varepsilon_j.
\end{equation}
Small pointwise error does not certify pairs whose conditional risks are nearly tied.
This is why a quadrature bound should not be presented as an AURC bound without an
additional margin argument.

\section{Dice instantiation and relation to SDC}\label{sec:theory}

For binary masks, define
\begin{equation}\label{eq:dice-definition}
  \Dice(A,B)=\frac{2|A\cap B|}{|A|+|B|},
\end{equation}
with $\Dice(\varnothing,\varnothing)=1$ and Dice equal to zero when exactly one mask
is empty. The loss $L_{\Dice}=1-\Dice$ lies in $[0,1]$. The framework gives
\begin{equation}\label{eq:levelset-dice}
\begin{split}
  C_{\Dice,M}(x)
  &:=-\sum_{m=1}^M w_m
      \left[1-\Dice(Y_{t_m},\Yhat)\right],\\
  1+C_{\Dice,M}(x)
  &=\sum_{m=1}^M w_m\Dice(Y_{t_m},\Yhat).
\end{split}
\end{equation}
Thus $1+C_{\Dice,M}$ is a quadrature estimate of expected Dice under $Q_p$ and is
rank-equivalent to the loss-oriented confidence.

For the ratio algebra below, assume that the deployed mask is nonempty, so the Dice
denominator is positive for every possible $Y$. Empty deployed masks remain covered
by the conventions above but require carrying the separately defined $0/0$ case
rather than treating Dice as an ordinary ratio.

SDC instead computes the soft Dice between $p$ and the deployed binary mask,
\begin{equation}\label{eq:sdc}
  \SDC(p,\Yhat)
  =\frac{2\sum_i p_i\hat y_i}
  {\sum_i p_i+\sum_i\hat y_i},
  \qquad
  \hat y_i=\mathbf 1\{i\in\Yhat\}.
\end{equation}
The degenerate zero-denominator case requires the same explicit empty convention as
Dice. In the nonempty setting just stated, SDC has the following exact interpretation.

\begin{proposition}[SDC is a ratio of expected Dice components]\label{prop:sdc}
Let $Q$ be any distribution over binary masks with pixel marginals
$\Ex_Q[Y_i]=p_i$. Define
\begin{equation*}
  N(Y)=2\sum_iY_i\hat y_i,
  \qquad
  D(Y)=\sum_iY_i+\sum_i\hat y_i.
\end{equation*}
Then
\begin{equation}\label{eq:sdc-ratio-expectations}
  \frac{\Ex_Q[N(Y)]}{\Ex_Q[D(Y)]}
  =\SDC(p,\Yhat).
\end{equation}
The value on the right is therefore invariant to the coupling among pixels once the
marginals are fixed.
\end{proposition}
\begin{proof}
Both $N$ and $D$ are affine in $Y$. Hence
$\Ex_Q[N]=2\sum_i p_i\hat y_i$ and
$\Ex_Q[D]=\sum_i p_i+\sum_i\hat y_i$, irrespective of higher-order dependence.
\end{proof}

Proposition~\ref{prop:sdc} does not make expected Dice coupling-invariant. Expected
Dice is an expectation of a ratio,
\begin{equation}\label{eq:expected-ratio}
  \Ex_Q[\Dice(Y,\Yhat)]
  =\Ex_Q\!\left[\frac{N(Y)}{D(Y)}\right],
\end{equation}
whereas SDC is the ratio in \eqref{eq:sdc-ratio-expectations}. In general,
\begin{equation}\label{eq:ratio-noncommute}
  \Ex_Q\!\left[\frac{N(Y)}{D(Y)}\right]
  \neq
  \frac{\Ex_Q[N(Y)]}{\Ex_Q[D(Y)]}.
\end{equation}

The distinction appears with only two pixels. Let
$p=(1/2,1/2)$ and $\Yhat=\{1\}$. If the two label pixels are independent, direct
enumeration gives
\begin{equation*}
  \Ex_{Q_{\mathrm{ind}}}[\Dice(Y,\Yhat)]=\frac{5}{12}.
\end{equation*}
Under the level-set coupling, $Y$ is either $\{1,2\}$ or $\varnothing$, each with
probability $1/2$, so
\begin{equation*}
  \Ex_{Q_p}[\Dice(Y,\Yhat)]=\frac13.
\end{equation*}
In both cases the marginals are $p$, and
$\SDC(p,\Yhat)=1/2$. The expected Dice under the two couplings and SDC are three
different quantities.

\paragraph{What each approximation assumes.}
Computing SDC requires no explicit joint posterior: equation~\eqref{eq:sdc} is a
function of $(p,\hat y)$ alone. Interpreting it as an approximation to the true ideal
expected Dice requires more. The pixel scores must represent the relevant marginals,
and the random denominator in \eqref{eq:expected-ratio} must concentrate sufficiently
for a ratio of expectations to approximate an expectation of a ratio. Conditional
pixel independence is a sufficient concentration regime used in the analysis of SDC
\cite{sdc}; it is not required merely to evaluate the SDC formula.

The level-set estimator makes a different trade. It evaluates expected Dice exactly
under the explicit comonotone working posterior $Q_p$ in the dense-quadrature limit.
Its connection to the true ideal confidence requires the loss-specific discrepancy
in \eqref{eq:loss-discrepancy} to be small. This is a statement about the joint law,
not only the marginals.

The assumptions are not ordered. If the true labels are conditionally independent
Bernoulli pixels with nondegenerate marginals, an independence-based concentration
argument may be appropriate while $Q_p$ can be far from the true posterior. If the
true label is generated by one shared threshold, then $Q_p$ is exact while conditional
independence fails except in degenerate cases. Neither example implies the other.

The correct comparison is therefore:
\begin{itemize}
\item SDC is Dice-specific, marginal-only, and computationally closed form. Its
ratio-of-expectations target can approximate ideal expected Dice under suitable
calibration and concentration.
\item The level-set construction is indexed by a general bounded loss and retains an
explicit joint working model. It estimates expectation of the chosen loss under that
model, with a separate posterior-discrepancy term.
\end{itemize}
Our framework generalizes the \emph{loss-indexed principle} beyond Dice; it neither
contains SDC as a literal special case nor has uniformly weaker probabilistic
assumptions.

\section{Limitations}\label{sec:limitations}

\paragraph{The working posterior is not identified by a probability map.}
The shared-threshold construction selects one coupling from many with the same pixel
marginals. Marginal calibration, even if exact, does not identify the expected value
of a nonlinear geometric loss. Validating the coupling requires repeated labels for
the same image or another source of information about annotation variability; a
single reference mask per image cannot separate marginal calibration error from
dependence-model error.

\paragraph{The loss convention is substantive.}
Dice is naturally bounded, but surface distances are undefined when one mask is empty.
Assigning the image diagonal, a fixed cap, or another application cost produces
different valid losses and can change the induced ranking. The convention must be
chosen from deployment semantics and declared before evaluation. Dropping undefined
cases changes the evaluated population and should be treated as a sensitivity
analysis rather than a silent implementation detail.

\paragraph{Quadrature control is not posterior control.}
The bound in \eqref{eq:decomp} separates these terms but does not make either small by
assertion. Total variation bounds are usually unavailable, the loss-specific
discrepancy must be supported empirically or stated as an assumption, and
$V(H_{L,x})/(2M)$ can be loose. A larger $M$ is guaranteed to improve the worst-case
midpoint approximation rate to the $Q_p$ integral, not to improve AURC monotonically
on a finite dataset.

\paragraph{Probability-map-only confidence misses information absent from the map.}
If two cases produce the same $(p,\Yhat)$, every deterministic score in this framework
assigns them the same confidence, even if external context makes one more likely to be
wrong. Confidently wrong, sharply delineated predictions remain a general failure
mode. Features, auxiliary failure predictors, ensembles, or distribution-shift
detectors can add information but lie outside the single-map post-hoc setting studied
here.

\paragraph{Scope.}
The formulation above is written for binary masks to isolate the relationship among
the deployed action, the loss, and the score. Extensions beyond this binary setting
require additional label-space and aggregation choices and are outside the present
scope. Finally, AURC evaluates ranking; it does not imply that the numerical value of
$C_{L,M}$ is calibrated as a probability or as an absolute forecast of loss.

\section{Conclusion}

Selective segmentation confidence is meaningful only relative to a deployed mask and
a deployment loss. We defined the corresponding ideal confidence, represented a
classical level-set construction as an explicit working posterior, and obtained a
single quadrature form for any bounded measurable segmentation loss. The resulting
theory makes two errors visible rather than conflating them: numerical approximation
of the working-posterior integral and disagreement between that posterior and the true
label law. Dice illustrates why the distinction matters. Expected Dice under the
level-set law and SDC's ratio of expected components are different estimands with
incomparable probabilistic assumptions. This loss-indexed view supports a direct
experimental question: how do confidence scores induced by different losses rank
predictions under the corresponding and cross-matched selective risks?
